\subsection{组合分析}

命题 9. 设 E 和 F 为两个集合，a 和 b 为它们的基数，f 为 E 到 F 上的一个满射，使得集合 \(\overline{f}^{1}(y)\) （其中 \(y \in F\) ）都具有相同的基数 c。则 a = bc。

因为族 \((\bar{f}^1 (y))_{\gamma \in \mathbb{F}}\) 是 \(\mathbf{E}\) 的一个划分，且划分的每个元素都是一个基数为 \(\mathfrak{c}\) 的集合；因此结果由（§ 3，第 4 号，命题 6，推论 2）得出。

定义 2. 设 \(n\) 为一个整数。乘积 \(\prod_{i < n} (i + 1)\) 记作 \(n!\) （读作“阶乘 \(n\) '').

我们有 \(0! = 1\) （第二章，§5，第 3 号）以及 \(1! = 1\) 。对于每个整数 \(n\) 我们拥有 \((n + 1)! = n!(n + 1)\) 。此关系式连同关系式 \(0! = 1\) 刻画了项 \(n!\) ，这很容易通过对 \(n\) 进行归纳而看出.

命题10。设 \(m\) 和 \(n\) 是满足 \(m \leqslant n\) 的整数。则 \(n! / (n - m)!\) 是从具有 \(m\) 个元素的集合 \(A\) 到具有 \(n\) 个元素的集合 \(B\) 的单射数目。

证明通过对 A 的元素个数 \(m \leqslant n\) 进行归纳。如果 \(m = 0\) ，结果是显然的。假设 \(m + 1 \leqslant n\) 。设 A 是一个具有 \(m + 1\) 个元素的集合，设 \(A'\) 是 A 的一个具有 \(m\) 个元素的子集，并令 \(\{a\} = A - A'\) 。设 F, \(F'\) 分别为 A 到 B 的、 \(A'\) 到 B 的单射映射所成的集合，并设 \(\varphi\) 为映射 \(f \to f|A'\) ，它将每个函数 \(f \in F\) 映到其在 \(A'\) 上的限制。对于每一个 \(f' \in F'\) ，一个元素 \(f\) （ \(\overline{\varphi}(f')\) 的元素）由它的值唯一确定，即 \(f(a)\) ；由于 \(f\) 是单射，我们必须有 \(f(a) \in B - f'(A')\) . 由此可知 \(\overline{\varphi}(f')\) 的元素个数 \(n - m\) 与 \(B - f'(A')\) 的相同. 因此，根据命题9，F具有

\[
(n - m) \frac {n !}{(n - m) !} = \frac {n !}{(n - m - 1) !}
\]

个元素，根据归纳假设。这就完成了证明。

推论。具有 n 个元素的有限集合的排列数等于 n!。

因为这个数等于该集合到自身的单射的个数（§ 4，第 2 号，命题 2，推论 4）。

命题 11。设 E 为具有 n 个元素的有限集合，并设 \((p_i)_{1 \leqslant i \leqslant h}\) 为满足 \(\sum_{i=1}^{h} p_i = n\) 的有限整数序列。则 E 的满足如下条件的覆盖 \((\mathbf{X}_i)_{1 \leqslant i \leqslant h}\) ——诸集合 \(\mathbf{X}_i\) 两两不相交且 \(\operatorname{Card}(\mathbf{X}_i) = p_i\) （对于 \(1 \leqslant i \leqslant h\) ）——的个数等于

\[
n! / \prod_ {i = 1} ^ {h} p _ {i}!.
\]

设 G 为 E 的排列的集合，设 P 为由满足命题条件的诸覆盖 \((\mathbf{X}_{i})_{1\leqslant i\leqslant h}\) 组成的集合。由于 \(\sum_{i=1}^{h}p_{i}=n\) ，P 非空。设 \((\mathbf{A}_{i})_{1\leqslant i\leqslant h}\) 为 P 的一个元素。对于每个置换 \(f\in G\) 族 \((f(\mathbf{A}_{i}))_{1\leqslant i\leqslant h}\) 再次属于 P；我们将其记为 \(\varphi(f)\) 。对于每个元素 \((\mathbf{X}_{i})_{1\leqslant i\leqslant h}\) ，让我们计算置换的数目 \(f\in G\) 使得 \(\varphi(f)=(\mathbf{X}_{i})\) 。我们将有 \(\varphi(f)=(\mathbf{X}_{i})\) 当且仅当对于每个指标 i 我们有 \(f(\mathbf{A}_{i})=\mathbf{X}_{i}\) . 因此所考虑的排列 f 的集合与诸集合（从 \(A_{i}\) 到 \(X_{i}\) 上的映射所成的集合）的乘积等势（第二章，§4，第 7 号，命题 8）；从而集合 \(\overline{\varphi}((\mathbf{X}_{i})_{1\leqslant i\leqslant h})\) 具有 \(\prod_{i=1}^{h}p_{i}!\) 个元素（命题 10 的推论）。由于 G 有 n! 个元素，结论由命题 9 得出。

推论1. 设A是一个含有n个元素的集合，且p是一个整数 \(\leqslant n\) . 那么，A 的具有 p 个元素的子集的数量是 \(\frac{n!}{p!(n-p)!}\) 。在命题11中取 \(h=2, p_{1}=p, p_{2}=n-p\) 。

含 \(p\) 个元素的子集（ \(n\) 个元素的集合中的子集，其中 \(p \leqslant n\) ）的个数记为 \(\binom{n}{p}\) ，称为指标为 \(n\) 与 \(p\) 的二项式系数。由关系式 \(\binom{n}{p} = \frac{n!}{p!(n - p)!}\) 立即得出 \(\binom{n}{p} = \binom{n}{n - p}\) .

这也是以下事实的推论：若 E 是有 n 个元素的集合，则 \(X \rightarrow E - X\) 是从 E 的含 p 个元素的子集集合到含 n − p 个元素的子集集合的双射。

我们置 \(\binom{n}{p} = 0\) （对于每一对满足 \(p > n\) 的自然整数）. 按照这一约定，含 \(p\) 个元素的子集（ \(n\) 个元素的集合中的）的个数是 \(\binom{n}{p}\) （对于每个自然整数 \(p\) ）.

推论 2. 设 E 和 F 分别是含 p 个和 n 个元素的全序有限集。则 E 到 F 的严格递增映射的个数为 \(\binom{n}{p}\) .

因为这样的映射是 E 到 F 的单射（§1，第 12 号，命题 11），并且由于 E 和 F 是良序的（§4，第 4 号，命题 3 的推论 1），对于 F 的每个含 p 个元素的子集 X，恰好存在一个 E 到 X 上的严格递增映射（§2，第 5 号，定理 3）。

命题 12. 对于每个整数 \(n\) ，我们有

\[
\sum_ {p} \binom{n}{p} = 2 ^ {n}.
\]

因为若E是一个含n个元素的集合，等式左边是E的子集个数。现在应用§3第5号中的命题12。

命题13。如果n和p是整数，则

\[
\binom {n + 1} {p + 1} = \binom {n} {p + 1} + \binom {n} {p}.
\]

设 E 是一个含 \(n + 1\) 个元素的集合，设 P 为 E 的所有子集的集合，其中包含 \(p + 1\) 个元素，设 a 为 E 的一个元素，并令

\[
\mathbf {E} ^ {\prime} = \mathbf {E} - \{a \}.
\]

令 \(P'\) （相应地， \(P''\) ）表示 E 的如下子集组成的集合：这些子集包含 \(p + 1\) 个元素且包含 a（相应地，不包含 a）。集合 \(P''\) 是含 \(p + 1\) 个元素的子集（ \(E'\) 中的）组成的集合，因此有 \(\binom{n}{p+1}\) 个元素。映射 \(X \to X \cap E'\) 是 \(P'\) 到由 p 个元素组成的子集（ \(E'\) 的子集）集合上的一个双射，因而 \(P'\) 有 \(\binom{n}{p}\) 个元素。现在结果由 P 是不相交集合 \(P'\) 与 \(P''\) 的并集这一事实得出.

命题13也可以借助公式 \(\binom{n}{p} = \frac{n!}{p!(n - p)!}\) （对于 \(p \leqslant n\) ）通过简单计算来证明.

命题14。设 \(n\) 是一个整数 \(> 0\) 。则数 \(a_{n}\) （相应地， \(b_{n}\) ）——整数有序对 \((i, j)\) （满足 \(1 \leqslant i \leqslant j \leqslant n\) ，相应地 \(1 \leqslant i < j \leqslant n\) ）的个数——是 \(\frac{1}{2} n(n + 1)\) ((相应地 \(\frac{1}{2} n(n - 1)\) ).

因为 \(b_{n}\) 是 \([1, n]\) 的二元子集的数目；因此

\[
b _ {n} = \frac {n !}{2 ! (n - 2) !} = \frac {1}{2} n (n - 1).
\]

由此可求得 \(a_{n}\) ：注意，有序对 \((i, j)\) （满足 \(1 \leqslant i \leqslant j \leqslant n\) ）所成集合，是有序对 \((i, j)\) （满足 \(1 \leqslant i \leqslant j < n\) ）所成集合与对 \((i, i)\) （其中 \(1 \leqslant i \leqslant n\) ）所成集合之并。因此 \(a_{n} = n + b_{n} = \frac{1}{2} n(n + 1)\) .

推论。对于每个整数 \(n > 0\) ，我们有

\[
\sum_ {i = 1} ^ {n} i = \frac {1}{2} n (n + 1).
\]

在整数有序对的集合 A ——有序对 \((i, j)\) 满足 \(1 \leqslant i \leqslant j \leqslant n\) ——中，令 \(A_{k}\) 表示对 \((i, k)\) （其中 \(1 \leqslant i \leqslant k\) ，对于任意整数 \(k \leqslant n\) ）所成的子集。则 \(A_{k}\) 有k个元素。但 \((\mathrm{A}_{k})_{1 \leqslant k \leqslant n}\) 是A的一个划分；因此结论成立。

命题15。设 \(n\) 与 \(h\) 为整数，并设 \(E\) 为具有 \(h\) 个元素的集合。则映射 \(u\) （从 \(E\) 到 \([0, n]\) 中的映射）——满足 \(\sum_{x \in E} u(x) \leqslant n\) （相应地， \(\sum_{x \in E} u(x) = n\) ，对于 \(h > 0\) ）——的个数是

\[
\binom {n + h} {h} \left(\text { 分别 } \binom {n + h - 1} {h - 1}\right).
\]

令 \(\mathrm{A}(h, n)\) （相应地， \(\mathrm{B}(h, n)\) ）表示映射 \(u\) （从 \(\mathbf{E}\) 到 \([0, n]\) ）的数目，满足 \(\sum_{x \in \mathbf{E}} u(x) \leqslant n\) （相应地， \(\sum_{x \in \mathbf{E}} u(x) = n\) ，当 \(h > 0\) ）。先证明 \(\mathrm{A}(h - 1, n) = \mathrm{B}(h, n)\) 。为此，令 \(\mathbf{E}'\) 为 \(\mathbf{E}\) 的一个具有 \(h - 1\) 个元素的子集，并令 \(\{a\} = \mathbf{E} - \mathbf{E}'\) 。如果 \(u\) 是从 \(\mathbf{E}\) 到 \([0, n]\) 且满足 \(\sum_{x \in \mathbf{E}} u(x) = n\) 的映射，那么它的限制 \(u'\) （到 \(\mathbf{E}'\) 上）满足 \(\sum_{x \in \mathbf{E}'} u'(x) \leqslant n\) ，并且 \(u(a) = n - \sum_{x \in \mathbf{E}'} u'(x)\) 。反之，每个映射 \(u'\) （从 \(\mathbf{E}'\) 到 \([0, n]\) ）若满足 \(\sum_{x \in \mathbf{E}'} u'(x) \leqslant n\) ，都唯一确定一个映射 \(u\) （从 \(\mathbf{E}\) 到 \([0, n]\) ），使得 \(u'\) 是它的限制，并且它满足 \(\sum_{x \in \mathbf{E}'} u(x) = n\) .

我们接下来注意到，如果 \(\sum_{x\in E}u(x)\leqslant n\) ，那么要么 \(\sum_{x\in E}u(x) = n\) ，要么 \(\sum_{x\in E}u(x)\leqslant n - 1\) ，且这两种可能性是互斥的。因此

\[
\mathrm{A} (h, n) = \mathrm{A} (h, n - 1) + \mathrm{B} (h, n) = \mathrm{A} (h, n - 1) + \mathrm{A} (h - 1, n).
\]

由于 \(\mathbf{A}(0,0) = 1 = \binom{0}{0}\) ，公式 \(\mathbf{A}(h,n) = \binom{n + h}{h}\) 由上述内容及命题13，通过对 \(n + h\) 进行归纳而得出.

* 单项式 \(\mathbf{X}_1^{\alpha_1}\mathbf{X}_2^{\alpha_2}\ldots \mathbf{X}_h^{\alpha_n}\) （在 \(h\) 个未定元中，总次数为 \(\leqslant n\) ）的个数显然等于映射 \(i\to \alpha_{i}\) （从 \([1,h]\) 到 \([0,n]\) 中的映射）——满足 \(\sum_{i = 0}^{h}\alpha_{i}\leqslant n,\) ——的个数，因此根据命题15等于 \((n + h)\) ；这个数也是 \(h + 1\) 个未定元中总次数为 \(n.\) 的单项式的个数. *

